\section{Data}\label{data}

\subsection{Source and Sample
Construction}\label{source-and-sample-construction}

All data are sourced from the Bloomberg Terminal via the Bloomberg
Desktop API (\texttt{blpapi}). I construct three panel datasets serving
distinct roles in the analysis:

\begin{enumerate}
\def\labelenumi{\arabic{enumi}.}
\item
  \textbf{Primary S\&P~500 panel.} Daily observations from January~1,
  2025 through March~25, 2026, comprising 160,487 firm-day observations
  across 503 S\&P~500 constituents. This panel contains the full set of
  treatment, outcome, and control variables used in the FE-OLS and
  DoubleML specifications.
\item
  \textbf{Extended S\&P~500 panel.} Daily observations from January~1,
  2016 through February~27, 2026, comprising approximately 1.19~million
  firm-day observations across the same 503 firms. This panel contains a
  reduced variable set (prices, volume, market capitalization, sentiment
  scores, and Fama-French factors) and is used exclusively for the
  temporal Fama-MacBeth analysis that tracks the sentiment coefficient
  from 2016 to 2026.
\item
  \textbf{Russell 3000 panel.} Daily observations from January~1, 2016
  through February~27, 2026, comprising approximately 1.97~million
  firm-day observations across the Russell~3000 constituents. This panel
  is used for the cap-group decomposition analysis and contains prices,
  volume, market capitalization, and sentiment scores. The Russell 3000
  analysis employs the same Fama-MacBeth and Panel OLS specifications as
  the S\&P~500 analysis, enabling a direct universe comparison.
\end{enumerate}

Non-trading days (weekends, holidays) are excluded by dropping
observations where both opening and closing prices are missing.
Bloomberg-specific missing value codes (\texttt{\#N/A\ N/A},
\texttt{\#N/A\ Field\ Not\ Applicable}) are recoded as \texttt{NA} prior
to any numerical transformation.

\subsection{Variable Definitions}\label{variable-definitions}

Table~\ref{tab:variables} defines all variables used in
the analysis. The S\&P~500 primary panel uses the full variable set; the
extended S\&P~500 and Russell 3000 panels use the subset marked with a
dagger (\dag). Those where the ``Bloomberg Field'' is blank, reflect
variables computed by combining values from the other data fields.

\begin{longtable}[]{@{}llll@{}}
\caption{Variable Definitions\label{tab:variables}}\\
\toprule
Variable & Bloomberg Field & Definition & Panels\tabularnewline
\midrule
\endhead
& & &\tabularnewline
\texttt{return} & --- &
\(\text{px\_close}_{\text{it}} - \text{px\_open}_{\text{it}}\)
(open-to-close, \$) & Primary\tabularnewline
\texttt{return\_oo} & --- & Open-to-open return (\%) & Extended,
Russell\tabularnewline
\texttt{return\_oo\_adj} & --- & FF4-factor residual of
\texttt{return\_oo} & Extended, Russell\tabularnewline
& & &\tabularnewline
\texttt{twitter\_sent} & \texttt{TWITTER\_SENTIMENT\_DAILY\_AVG} & Daily
avg.~Twitter sentiment (\(\lbrack - 1,1\rbrack\)) & All\tabularnewline
\texttt{twitter\_neg\_count} & \texttt{TWITTER\_NEG\_SENTIMENT\_COUNT} &
Count of negative-polarity tweets & All\tabularnewline
\texttt{news\_sent} & \texttt{NEWS\_SENTIMENT\_DAILY\_AVG} & Daily
avg.~news sentiment (\(\lbrack - 1,1\rbrack\)) & All\tabularnewline
& & &\tabularnewline
\texttt{px\_high} & \texttt{PX\_HIGH} & Intraday high price (\$) &
Primary\tabularnewline
\texttt{px\_low} & \texttt{PX\_LOW} & Intraday low price (\$) &
Primary\tabularnewline
\texttt{mkt\_cap} & \texttt{CUR\_MKT\_CAP} & Market cap (millions, \$) &
All\tabularnewline
\texttt{log\_mkt\_cap} & --- & \(\text{log}(\text{mkt\_cap})\) &
Extended, Russell\tabularnewline
\texttt{total\_equity} & \texttt{TOTAL\_EQUITY} & Book equity (millions,
\$) & Primary\tabularnewline
\texttt{debt\_to\_equity} & \texttt{TOT\_DEBT\_TO\_TOT\_EQY} & Total
debt / total equity (\%) & Primary\tabularnewline
\texttt{volume} & \texttt{PX\_VOLUME} & Daily share volume (shares) &
All\tabularnewline
\texttt{abnorm\_vol} & --- & \% deviation of volume from ticker mean &
Extended, Russell\tabularnewline
\texttt{vol\_rs} & --- & Rogers-Satchell realized volatility & Extended,
Russell\tabularnewline
\texttt{rsi\_30} & \texttt{RSI\_30D} & 30-day Relative Strength Index &
Primary\tabularnewline
\texttt{ma\_50} & \texttt{MOV\_AVG\_50D} & 50-day moving average of
price (\$) & Primary\tabularnewline
\texttt{bid\_ask\_spread} & \texttt{AVERAGE\_BID\_ASK\_SPREAD} & Daily
avg.~bid-ask spread (\$) & All\tabularnewline
& & &\tabularnewline
\texttt{lagk} & --- & \(\text{return}_{i,t - k}\),
\(k \in \{ 1,2,3,5,7\}\) & Primary\tabularnewline
\texttt{cap\_group} & --- & Large / Mid / Small (see
Section~\ref{cap-group-classification}) &
Russell\tabularnewline
& & &\tabularnewline
\bottomrule
\end{longtable}

The primary outcome for the S\&P~500 specifications, \texttt{return}, is
the raw dollar price change from market open to close on day \(t\) for
firm \(i\). This differs from the percentage return used in Gu and Kurov (2020); the dollar
denomination preserves interpretability of treatment coefficients as
dollar-per-unit sentiment effects, and firm fixed effects absorb
cross-sectional differences in share price levels. For the extended and
Russell panels, the outcome is the factor-adjusted open-to-open return
(\texttt{return\_oo\_adj}), which is the appropriate dependent variable
for the Fama-MacBeth specifications following Gu and Kurov (2020).

The primary treatment variable, \texttt{twitter\_sent}, is Bloomberg's
proprietary daily average Twitter sentiment score. Bloomberg applies
natural language processing to the universe of public tweets mentioning
a given firm's name, ticker, or cashtag, assigning each tweet a score on
\(\lbrack - 1, + 1\rbrack\) and averaging across all tweets within the
calendar day. The variable \texttt{twitter\_neg\_count} counts only
those tweets classified with negative polarity, following Teti et al.\ (2019). The news
sentiment variable, \texttt{news\_sent}, is constructed analogously from
Bloomberg-sourced news articles.

\subsection{Cap-Group Classification}\label{cap-group-classification}

For the Russell 3000 analysis, firms are assigned to one of three
market-capitalization groups based on their time-averaged market
capitalization over the full sample period:

\begin{itemize}
\item
  \textbf{Large cap:} \({\bar{\text{mkt\_cap}}}_{i} > \$ 10\)B
  (\(N = 763,274\) firm-days)
\item
  \textbf{Mid cap:} \(\$ 2\)B
  \(\leq {\bar{\text{mkt\_cap}}}_{i} \leq \$ 10\)B (\(N = 640,932\)
  firm-days)
\item
  \textbf{Small cap:} \({\bar{\text{mkt\_cap}}}_{i} < \$ 2\)B
  (\(N = 568,068\) firm-days)
\end{itemize}

Group membership is fixed for the duration of the sample using the
time-average rather than daily market capitalization, so that firms do
not migrate between groups due to short-run price volatility. This
ensures that within-group temporal comparisons are not confounded by
within-firm capitalization shifts.

The large-cap group corresponds to an expanded version of the S\&P~500
universe, with the mid-cap and small-cap groups to the Russell 2000. The
thresholds of \$2B and \$10B are standard in institutional investment
practice and are consistent with the Bloomberg index classification
methodology.

\subsection{Measurement Limitations}\label{measurement-limitations}

Three features of the Bloomberg sentiment data warrant discussion.

\textbf{Proprietary classification.} Bloomberg does not disclose the
dictionary, model architecture, or training data underlying its
sentiment classification. This precludes replication with alternative
NLP tools and prevents diagnostic checks on whether specific tweet
categories (e.g., bot-generated, reply threads, retweets) are included
or filtered. The Loughran and McDonald (2011) critique of general-purpose sentiment dictionaries
applies with unknown force: it is possible that Bloomberg's classifier
misattributes domain-specific financial language.

\textbf{Daily Aggregation.} The sentiment scores are daily averages.
Bloomberg aggregates all tweets within a 24-hour window without
distinguishing pre-market, intraday, and after-hours activity. Because
the primary outcome variable is measured from market open to market
close, a tweet posted after the close is counted in the same day's
sentiment as one posted before the open. This aggregation mismatch
biases the design against detecting clean causal effects and may
introduce reverse causality if after-hours tweets respond to the day's
price action that has already been recorded in \texttt{return}. This
concern is directly assessed by the placebo test in
Section~\ref{robustness-checks}.

\textbf{Selection in sentiment coverage.} Not all firm-days have Twitter
sentiment scores. Missing sentiment days are not randomly distributed
across firms or time: large-cap, high-volume firms receive
disproportionate coverage. In the S\&P~500 extended panel, Twitter
sentiment is non-missing for 908,781 of 1,190,286 firm-days (76\%
coverage). Coverage is sparser for Russell~3000 firms, particularly for
the small-cap segment where lower public visibility reduces the daily
tweet count below Bloomberg's classification threshold on many days.
Summary statistics for all three panels are reported in
Tables~\ref{tab:sumstats_primary}--\ref{tab:sumstats_russel} and document
the extent of missingness across key variables.

\subsection{Sample Properties}\label{sample-properties}

\textbf{S\&P~500 primary panel.} The panel spans 503 firms over
approximately 320 trading days (January~2025--March~2026), yielding
160,487 firm-day observations with non-missing price data. The effective
sample size for the full FE-OLS specification (Model~2) drops to 147,128
observations after listwise deletion of rows with missing controls,
primarily \texttt{total\_equity} and \texttt{debt\_to\_equity} for newly
listed firms. The baseline specification (Model~1) retains 160,103
observations. For DoubleML, XGBoost handles missing covariates natively;
observations are excluded only when the treatment variable itself is
missing. Summary statistics are reported in
Table~\ref{tab:sumstats_primary}. Twitter sentiment is non-missing for
160,103 of 160,487 firm-days (99.8\% coverage); the high coverage
reflects the S\&P~500's large-cap composition, which generates sufficient
daily tweet volume to meet Bloomberg's classification threshold on nearly
every trading day.

\textbf{Extended S\&P~500 panel.} The panel covers 2,633 unique trading
days across the same 503 firms (January~2016--February~2026), yielding
approximately 1.19~million firm-day observations. Twitter sentiment
coverage begins in earnest around 2017 for most firms, with sporadic
coverage in 2016. The Fama-MacBeth regressions exclude dates on which
fewer than 30 firms report non-missing sentiment. The sample is truncated
at February~28, 2026, the last available date in the Kenneth French
factor library, to ensure clean four-factor risk adjustment throughout.
Twitter sentiment is non-missing for 908,781 of 1,190,286 firm-days
(76\% coverage). Summary statistics are reported in
Table~\ref{tab:sumstats}.

\textbf{Russell 3000 panel.} The panel covers approximately 2,921 firms
across the 2016--2026 window, yielding 1,972,274 firm-day observations
with non-missing price data. After restricting to firm-days with
non-missing \texttt{twitter\_sent}, the effective sample for Fama-MacBeth
estimation is smaller (2,183,113 firm-days with non-missing sentiment
across the full raw file); Fama-MacBeth regressions exclude dates on
which fewer than 30 firms per cap group report non-missing sentiment. The
large-cap group (\(N = 763,274\)) has the highest sentiment coverage
given its overlap with the S\&P~500 universe; the small-cap group
(\(N = 568,068\)) has the lowest, as lower public visibility reduces
daily tweet counts below Bloomberg's classification threshold on many
days. Summary statistics are reported in Table~\ref{tab:sumstats_russel}.
